% English translation of checked_pages/esq_119_038.tex; source: cote 119, site 38, PDF page 39, typescript page 36.
% AI-translated by GPT-6 Astra (OpenAI, ChatGPT Pro mode).
\scanpage{119/38}
\mspage{36}

involved in the situation, the theory of dévissage that
I had reached\frapparatus{V-FR-038-001} in an entirely different context --- with the
additional task, subsequently, of drawing out a general comparison
theorem, on the model of the well-known theorems, to
compare the invariants obtained (in particular the homotopy types of
various tubular neighborhoods) in the transcendental setting and in
the scheme-theoretic setting. I was able to convince myself that such a formalism
of dévissage did indeed make sense in the (so-called
``abstract'' !) context of general topoi, or at least of noetherian
topoi (such as those that arise here), through a suitable
notion of the \uline{canonical tubular neighborhood of a subtopos} in an
ambient topos. Once this notion has been acquired, with certain simple
formal properties, the description of the ``dévissage'' of a
stratified topos is considerably simpler even in this
setting than in the topological (tame) setting. It is true that
here too work on foundations remains to be done, in particular on
the very notion of the tubular neighborhood of a subtopos --- and it is
surprising, moreover, that this work (as far as I know)
has still not been done, that is, that no one (in the more
than twenty years that a context of étale topologies has existed) seems
to have needed it; surely a sign that the understanding
of the topological structure of schemes has not progressed all that
much since the work of Artin-Mazur ...

Once the double task of preliminary work
(more or less heuristic) around the notion of the dévissage of a
stratified space or topos had been accomplished, which was a crucial stage in
my understanding of moduli multiplicities, it moreover
became apparent that for the needs of the latter, one can probably
short-circuit at least a good part of this theory by
direct geometric arguments. The fact remains that for
me, the formalism of dévissage that I reached has proved
its usefulness and coherence, independently of any
question about the most appropriate foundations that would allow
it to be given its full meaning.

\margin{38}
